\documentclass[12pt]{article}                                                                                                           
\usepackage[utf8]{inputenc}                                                                                                             
\usepackage[T1]{fontenc}                                                                                                                
\usepackage{amsmath, amssymb}                                                                                                           
\usepackage{geometry}                                                                                                                   
\usepackage{xcolor}                                                                                                                     
\usepackage{lipsum}                                                                                                                     
\usepackage{titlesec}                                                                                                                   
\usepackage{fancyhdr}                                                                                                                   
\usepackage[most]{tcolorbox}                                                                                                                  
\usepackage{graphicx}                                                                                                                   
\usepackage{float}                                                                                      
\usepackage{tikz}
\usetikzlibrary{shapes.geometric, arrows.meta}
\usetikzlibrary{arrows.meta, decorations.markings}                                                                                      
                                                                                                                                        
\geometry{a4paper, margin=2.5cm}                                                                                                        
\pagestyle{fancy}                                                                                                                       
\fancyhf{}                                                                                                                              
\rhead{Problem Set I}                                                                                                                   
\lhead{Algebra IV}                                                                                                                     
\cfoot{\thepage}                                                                                                                        
                                                                                                                                        
% Fancy titles                                                                                                                          
\titleformat{\section}{\normalfont\Large\bfseries}{Exercise \thesection:}{1em}{}                                                        
\titleformat{\subsection}{\normalfont\bfseries}{Correction:}{1em}{}                                                                     
                                                                                                                                        
% Custom box for correction with breakable option
\newtcolorbox{correctionbox}{                                                                                                           
  colback=gray!5,                                                                                                                       
  colframe=black,                                                                                                                       
  fonttitle=\bfseries,                                                                                                                  
  title=Correction,                                                                                                                     
  before skip=10pt,                                                                                                                     
  after skip=10pt,
  breakable,
  enhanced
}                                                                                                                                       
                                                                                                                                        
% Fix headheight warning
\setlength{\headheight}{14.49998pt}

\begin{document}

\begin{center}
\Large\textbf{Ibn Tofail University} \\[1em]
\large\textit{Algebra IV} \\[2em]
\large\textit{Problem Set I} \\[0.5em]
\end{center}

\vspace{1cm}

% ----------- EXERCISE 1 -------------
\section{}
Let \( f \in \text{End}(\mathbb{C}^3) \) with matrix relative to the canonical basis:
\[
A = \begin{pmatrix}
2 & 0 & 4 \\
3 & -4 & 12 \\
1 & -2 & 5
\end{pmatrix}
\]

\begin{enumerate}
    \item Find the eigenvalues of \( f \).
    \item Show there exists a basis \( B \) of eigenvectors of \( f \) and give the matrix of \( f \) in \( B \).
    \item Give the change of basis matrix \( P \) from the canonical basis \( B_0 \) to \( B \), then compute \( P^{-1}AP \).
\end{enumerate}

\begin{correctionbox}
    \begin{enumerate}
        \item Finding the eigenvalues
\begin{align*}
\det(A - \lambda I) = \lambda(2-\lambda)(\lambda - 1)
\end{align*}

    \item Finding eigenvectors and basis \(B\)

\textbf{For \(\lambda = 0\):} Solve \(A\mathbf{v} = 0\)
\[
\Rightarrow x = -2z,\ y = \frac{3}{2}z
\]
Take \(z = 2\): \(v_1 = (-4, 3, 2)\)

\textbf{For \(\lambda = 1\):} Solve \((A - I)\mathbf{v} = 0\)
\[
\Rightarrow x = -4z,\ y = 0
\]
Take \(z = 1\): \(v_2 = (-4, 0, 1)\)

\textbf{For \(\lambda = 2\):} Solve \((A - 2I)\mathbf{v} = 0\)
\[
\Rightarrow z = 0,\ x = 2y
\]
Take \(y = 1\): \(v_3 = (2, 1, 0)\)

Basis \(B = \{v_1, v_2, v_3\} = \{(-4,3,2), (-4,0,1), (2,1,0)\}\)

Matrix of \(f\) in basis \(B\):
\[
D = \begin{pmatrix}
0 & 0 & 0 \\
0 & 1 & 0 \\
0 & 0 & 2
\end{pmatrix}
\]

    \item Change of basis matrix and similarity transformation

Change of basis matrix:
\[
P = \begin{pmatrix}
-4 & -4 & 2 \\
3 & 0 & 1 \\
2 & 1 & 0
\end{pmatrix}
\]

We verify that:
\[
P^{-1}AP = D = \begin{pmatrix}
0 & 0 & 0 \\
0 & 1 & 0 \\
0 & 0 & 2
\end{pmatrix}
\]

This confirms that \(A\) is diagonalizable and \(B\) is a basis of eigenvectors.
    \end{enumerate}
\end{correctionbox}

% ----------- EXERCISE 2 -------------
\section{}
Let 
\[
A = \begin{pmatrix}
2 & -1 & 1 \\
2 & 0 & 1 \\
1 & 1 & 1
\end{pmatrix}
\]
Find the characteristic polynomial of \( A \).

\begin{correctionbox}
The characteristic polynomial is given by:
\[
P(\lambda) = \det(A - \lambda I)
\]
where \(I\) is the \(3 \times 3\) identity matrix.

We compute, Using the cofactor expansion along the first row:
\begin{align*}
\det(A - \lambda I) &= (2-\lambda)\left[(-\lambda)(1-\lambda) - (1)(1)\right] \\
&\quad - (-1)\left[(2)(1-\lambda) - (1)(1)\right] \\
&\quad + (1)\left[(2)(1) - (-\lambda)(1)\right]\\
\iff P(\lambda) &= (2-\lambda)(\lambda^2 - \lambda - 1) + (1 - 2\lambda) + (2 + \lambda) \\
&= -\lambda^3 + 3\lambda^2 - 2\lambda + 1
\end{align*}
\end{correctionbox}

% ----------- EXERCISE 3 -------------
\section{}
Let \( A \in M_K \) with characteristic polynomial \( P \) such that \( P(0) \neq 0 \).
\begin{itemize}
    \item Show \( A \) is invertible and give the expression for the characteristic polynomial of \( A \).
    \item Show for all \( B \in M_K \), \( AB \) and \( BA \) have the same characteristic polynomial.
\end{itemize}

\begin{correctionbox}
\begin{enumerate}
    \item Since $ P(\lambda) = \det(A - \lambda I) $, we have $ P(0) = \det(A) $. The hypothesis $ P(0) \neq 0 $ implies $ \det(A) \neq 0 $, so $ A $ is invertible. 
    
    Let $ Q(\lambda) $ be the characteristic polynomial of $ A^{-1} $. Then:
    \[
    Q(\lambda) = \det(A^{-1} - \lambda I) = \det(A^{-1}(I - \lambda A)) = \det(A^{-1})\det(I - \lambda A).
    \]
    Since $ \det(A^{-1}) = 1/\det(A) = 1/P(0) $ and 
    \[
    \det(I - \lambda A) = \det(-\lambda(A - \lambda^{-1}I)) = (-\lambda)^n \det(A - \lambda^{-1}I) = (-\lambda)^n P(\lambda^{-1}),
    \]
    we obtain
    \[
    Q(\lambda) = \frac{1}{P(0)} \cdot (-\lambda)^n P(\lambda^{-1}) = (-1)^n \frac{\lambda^n}{P(0)} P\left(\frac{1}{\lambda}\right).
    \]

    \item Since $ A $ is invertible (from the first part), for any $ B \in M_n(K) $ we have
    \[
    BA = A^{-1}(AB)A,
    \]
    which shows that $ AB $ and $ BA $ are similar matrices. Similar matrices have the same characteristic polynomial, so $ AB $ and $ BA $ share the same characteristic polynomial.
\end{enumerate}
\end{correctionbox}

% ----------- EXERCISE 4 -------------
\section{}
Given 
\[
A = \begin{pmatrix}
3 & 1 & 1 \\
2 & 4 & 2 \\
1 & 1 & 3
\end{pmatrix}
\]

\begin{enumerate}
    \item Determine the characteristic polynomial of \( A \).
    \item Determine eigenvalues and give a basis for each eigenspace of \( A \).
    \item Is \( A \) triangulable? Diagonalizable?
\end{enumerate}

\begin{correctionbox}
\begin{enumerate}
    \item The characteristic polynomial of $A$ is $P(\lambda) = \det(A - \lambda I)$.
    \begin{align*}
    P(\lambda) &= (3-\lambda)\det\begin{pmatrix}4-\lambda & 2 \\ 1 & 3-\lambda\end{pmatrix} 
    - 1\cdot\det\begin{pmatrix}2 & 2 \\ 1 & 3-\lambda\end{pmatrix} 
    + 1\cdot\det\begin{pmatrix}2 & 4-\lambda \\ 1 & 1\end{pmatrix} \\
    &= -(\lambda-2)^2(\lambda-6).
    \end{align*}
    
    \item The eigenvalues are the roots of $P(\lambda)$: $\lambda_1 = 2$ (algebraic multiplicity 2) and $\lambda_2 = 6$ (algebraic multiplicity 1).
    \textbf{Eigenspace for $\lambda = 2$:} Solve $(A - 2I)v = 0$:
    \[
    A - 2I = \begin{pmatrix}
    1 & 1 & 1 \\
    2 & 2 & 2 \\
    1 & 1 & 1
    \end{pmatrix} \xrightarrow{\text{row reduction}} \begin{pmatrix}
    1 & 1 & 1 \\
    0 & 0 & 0 \\
    0 & 0 & 0
    \end{pmatrix}.
    \]
    The equation is $x + y + z = 0$, so $x = -y - z$. A basis is:
    \[
    \left\{ \begin{pmatrix}-1 \\ 1 \\ 0\end{pmatrix}, \begin{pmatrix}-1 \\ 0 \\ 1\end{pmatrix} \right\}.
    \]

    \textbf{Eigenspace for $\lambda = 6$:} Solve $(A - 6I)v = 0$:
    \[
    A - 6I = \begin{pmatrix}
    -3 & 1 & 1 \\
    2 & -2 & 2 \\
    1 & 1 & -3
    \end{pmatrix} \xrightarrow{\text{row reduction}} \begin{pmatrix}
    1 & 0 & -1 \\
    0 & 1 & -2 \\
    0 & 0 & 0
    \end{pmatrix}.
    \]
    The equations are $x = z$ and $y = 2z$. A basis is:
    \[
    \left\{ \begin{pmatrix}1 \\ 2 \\ 1\end{pmatrix} \right\}.
    \]
    The geometric multiplicity is 1.

    \item Since all eigenvalues are real (assuming $K = \mathbb{R}$ or $\mathbb{C}$) and the characteristic polynomial splits into linear factors, $A$ is triangulable over $K$.
\end{enumerate}
\end{correctionbox}

% ----------- EXERCISE 5 -------------
\section{}
Let \( n \in \mathbb{N} \), \( K \) a commutative field, and \( A \in M \).
\begin{itemize}
    \item If \( \mathrm{rank}(A) = 1 \), show \( A \) is diagonalizable if and only if \( \mathrm{tr}(A) = 0 \).
\end{itemize}

\begin{correctionbox}
Let $ n \in \mathbb{N} $, $ K $ a commutative field, and $ A \in M_n(K) $ with $ \mathrm{rank}(A) = 1 $.

Since $ \mathrm{rank}(A) = 1 $, the image of $ A $ is a one-dimensional subspace of $ K^n $. Therefore, there exist non-zero vectors $ u, v \in K^n $ such that 
\[
A = u v^\top,
\]
where $ v^\top $ denotes the transpose of $ v $ (so $ A_{ij} = u_i v_j $).

The characteristic polynomial of a rank-1 matrix has a well-known form. Indeed, for any $ \lambda \in K $,
\[
\det(A - \lambda I) = (-\lambda)^{n-1} \bigl( \mathrm{tr}(A) - \lambda \bigr).
\]
This can be seen by noting that $ 0 $ is an eigenvalue of $ A $ with geometric multiplicity at least $ n - 1 $ (since $ \dim \ker A = n - \mathrm{rank}(A) = n - 1 $), and the remaining eigenvalue must be $ \mathrm{tr}(A) $, because the trace equals the sum of all eigenvalues (counted with algebraic multiplicities).

Thus, the eigenvalues of $ A $ are:
\[
\lambda_1 = \mathrm{tr}(A) \quad \text{(algebraic multiplicity 1)}, \qquad \lambda_2 = 0 \quad \text{(algebraic multiplicity } n-1\text{)}.
\]

Now, $ A $ is diagonalizable if and only if, for each eigenvalue, its geometric multiplicity equals its algebraic multiplicity.

- The eigenspace for $ \lambda = 0 $ is $ \ker A $, which has dimension $ n - \mathrm{rank}(A) = n - 1 $. Hence, the geometric multiplicity of $ 0 $ is $ n - 1 $, which already matches its algebraic multiplicity.

- For $ \lambda = \mathrm{tr}(A) $, the geometric multiplicity is $ 1 $ if $ \mathrm{tr}(A) \neq 0 $ (since the image of $ A $ is one-dimensional and consists of eigenvectors with eigenvalue $ \mathrm{tr}(A) $), and $ 0 $ if $ \mathrm{tr}(A) = 0 $ (because then $ \mathrm{tr}(A) = 0 $ is not a distinct eigenvalue).

We now consider two cases:

\textbf{Case 1: $ \mathrm{tr}(A) \neq 0 $.}  
Then the eigenvalues are $ \mathrm{tr}(A) $ (multiplicity 1) and $ 0 $ (multiplicity $ n-1 $). The geometric multiplicity of $ \mathrm{tr}(A) $ is 1 (since $ \mathrm{Im}(A) \subseteq E_{\mathrm{tr}(A)} $ and $ \dim \mathrm{Im}(A) = 1 $), so all geometric multiplicities match algebraic multiplicities. Hence, $ A $ is diagonalizable.

\textbf{Case 2: $ \mathrm{tr}(A) = 0 $.}  
Then the only eigenvalue is $ 0 $, with algebraic multiplicity $ n $. However, the geometric multiplicity of $ 0 $ is $ \dim \ker A = n - 1 < n $. Therefore, $ A $ is \textbf{not} diagonalizable.

But wait—this contradicts the statement we are asked to prove. The issue is in the interpretation: the correct statement is actually:

\[
\text{If } \mathrm{rank}(A) = 1, \text{ then } A \text{ is diagonalizable } \iff \mathrm{tr}(A) \neq 0.
\]

However, the problem claims the equivalence with $ \mathrm{tr}(A) = 0 $. This suggests either a typo in the problem or a different convention.

Let us double-check with an example:

- Take $ A = \begin{pmatrix} 1 & 0 \\ 0 & 0 \end{pmatrix} $. Then $ \mathrm{rank}(A) = 1 $, $ \mathrm{tr}(A) = 1 \neq 0 $, and $ A $ is diagonal (hence diagonalizable).

- Take $ A = \begin{pmatrix} 0 & 1 \\ 0 & 0 \end{pmatrix} $. Then $ \mathrm{rank}(A) = 1 $, $ \mathrm{tr}(A) = 0 $, but $ A $ is a non-zero nilpotent matrix, so it is not diagonalizable.

Thus, the correct equivalence is:  
\[
A \text{ is diagonalizable } \iff \mathrm{tr}(A) \neq 0.
\]

Therefore, the statement in the problem appears to contain an error. Assuming the intended statement is the correct one, we conclude:

\[
\boxed{A \text{ is diagonalizable if and only if } \mathrm{tr}(A) \neq 0.}
\]

If, however, the problem insists on the given formulation, then it is false as stated.
\end{correctionbox}

% ----------- EXERCISE 6 -------------
\section{}
Show matrix \( A \in M_K \) with entries \( a_{ij} = 1 \) for all \( i,j \) is diagonalizable.\\
Determine eigenvalues and eigenspaces of the linear map defined by
\[
f(X) = X \cdot (P X_1 X - 3P - XP)
\]

\begin{correctionbox}
We address the two parts of the problem separately.

\medskip

\noindent\textbf{Part 1: Diagonalizability of the all-ones matrix.}

Let $ A \in M_n(K) $ be the matrix defined by $ a_{ij} = 1 $ for all $ i,j $. This is the rank-1 matrix
\[
A = \mathbf{1}\,\mathbf{1}^\top,
\]
where $ \mathbf{1} = (1,1,\dots,1)^\top \in K^n $.

Since $ \operatorname{rank}(A) = 1 $ (all rows are identical and non-zero), we know from standard results that:
\begin{enumerate}
    \item $ 0 $ is an eigenvalue with geometric multiplicity $ n - 1 $ (because $ \dim \ker A = n - \operatorname{rank}(A) = n - 1 $).
    \item The trace of $ A $ is $ \operatorname{tr}(A) = \sum_{i=1}^n a_{ii} = n $.
\end{enumerate}

For a rank-1 matrix, the non-zero eigenvalue (if any) equals the trace. Hence the eigenvalues are:
\[
\lambda_1 = n \quad \text{(algebraic multiplicity 1)}, \qquad \lambda_2 = 0 \quad \text{(algebraic multiplicity } n-1\text{)}.
\]

Now, the eigenspace for $ \lambda = n $ is $ \operatorname{Im}(A) = \operatorname{span}\{\mathbf{1}\} $, which is 1-dimensional.  
The eigenspace for $ \lambda = 0 $ is $ \ker A = \{ x \in K^n \mid x_1 + x_2 + \cdots + x_n = 0 \} $, which has dimension $ n - 1 $.

Thus, the sum of the geometric multiplicities is $ 1 + (n - 1) = n $, so $ A $ is diagonalizable (provided that $ n \neq 0 $ in $ K $; i.e., the characteristic of $ K $ does not divide $ n $). In particular, over $ \mathbb{R} $ or $ \mathbb{C} $, $ A $ is always diagonalizable.

\medskip

\noindent\textbf{Part 2: Eigenvalues and eigenspaces of } $ f(X) = X \cdot (P X_1 X - 3P - XP) $.

The expression for the linear map $ f $ is ambiguous and likely contains typographical errors. The notation $ X \cdot (P X_1 X - 3P - XP) $ is not standard for a linear map on a space of matrices or vectors:
\begin{enumerate}
    \item It is unclear what $ X_1 $ denotes (a fixed matrix? a coordinate?).
    \item The term $ P X_1 X $ suggests a product of three objects, but without knowing the dimensions or nature of $ P $, $ X_1 $, and $ X $, the expression is undefined.
    \item Moreover, the map as written is generally \textbf{not linear} in $ X $, due to the presence of quadratic terms like $ X P X $ or $ P X_1 X $.
\end{enumerate}
Since the problem states that $ f $ is a linear map, the given formula must be incorrect or miswritten. A common linear map involving a fixed matrix $ P $ is, for example:
\[
f(X) = PX - XP \quad \text{(the commutator map)},
\]
or
\[
f(X) = PX + XP,
\]
both of which are linear in $ X $.

Without a clear and mathematically consistent definition of $ f $, it is impossible to determine its eigenvalues and eigenspaces.

\medskip

\noindent\textbf{Conclusion:}
\begin{enumerate}
    \item The all-ones matrix is diagonalizable (over fields where $ n \neq 0 $), with eigenvalues $ n $ (multiplicity 1) and $ 0 $ (multiplicity $ n-1 $).
    \item The second part cannot be solved as stated due to ambiguous or erroneous notation.
\end{enumerate}
\end{correctionbox}


% ----------- EXERCISE 7 -------------
\section{}
For \( n \in \mathbb{N} \) and two \( n \times n \) matrices \( A, B \), show every eigenvalue of \( AB \) is an eigenvalue of \( BA \).

\begin{correctionbox}
Let $ n \in \mathbb{N} $ and let $ A, B \in M_n(K) $. We show that every eigenvalue of $ AB $ is also an eigenvalue of $ BA $.

Let $ \lambda \in K $ be an eigenvalue of $ AB $. Then there exists a non-zero vector $ v \in K^n $ such that
\[
ABv = \lambda v.
\]

We consider two cases.

\medskip

\noindent\textbf{Case 1: $ \lambda \neq 0 $.}

Since $ \lambda \neq 0 $ and $ v \neq 0 $, we have $ Bv \neq 0 $ (otherwise $ ABv = A0 = 0 = \lambda v $ would imply $ \lambda v = 0 $, contradicting $ \lambda \neq 0 $ and $ v \neq 0 $).

Now apply $ B $ to both sides of $ ABv = \lambda v $:
\[
B(ABv) = B(\lambda v) \quad \Longrightarrow \quad (BA)(Bv) = \lambda (Bv).
\]
Since $ Bv \neq 0 $, this shows that $ Bv $ is an eigenvector of $ BA $ with eigenvalue $ \lambda $. Hence $ \lambda $ is an eigenvalue of $ BA $.

\medskip

\noindent\textbf{Case 2: $ \lambda = 0 $.}

We must show that if $ 0 $ is an eigenvalue of $ AB $, then it is also an eigenvalue of $ BA $. This is equivalent to showing that if $ AB $ is singular, then $ BA $ is singular.

Recall that for square matrices, $ \det(AB) = \det(A)\det(B) = \det(BA) $. Therefore,
\[
\det(AB) = 0 \quad \Longleftrightarrow \quad \det(BA) = 0.
\]
Thus, $ AB $ is not invertible if and only if $ BA $ is not invertible, which means $ 0 $ is an eigenvalue of $ AB $ if and only if $ 0 $ is an eigenvalue of $ BA $.

\medskip

In both cases, every eigenvalue $ \lambda $ of $ AB $ is also an eigenvalue of $ BA $. (In fact, the same argument with $ A $ and $ B $ swapped shows the converse, so $ AB $ and $ BA $ have exactly the same eigenvalues, including algebraic multiplicities when $ K $ is algebraically closed.)

Therefore, every eigenvalue of $ AB $ is an eigenvalue of $ BA $.
\end{correctionbox}

% ----------- EXERCISE 8 -------------
\section{}
For what values of real parameters \( a,b,c,d,e,f \) are the following matrices diagonalizable?
\[
1.
\quad A = \begin{pmatrix}
a & b & c \\
0 & 2 & d \\
e & 0 & 2
\end{pmatrix}
\quad
2.
\quad B = \begin{pmatrix}
a & b & c \\
0 & 1 & d \\
e & 0 & 2
\end{pmatrix}
\]

\begin{correctionbox}
The matrix 
\[
A = \begin{pmatrix}
a & b & c \\
0 & 2 & d \\
e & 0 & 2
\end{pmatrix}
\]
is diagonalizable over $\mathbb{R}$ **if and only if**:

\[
e = 0 \quad \text{and} \quad
\begin{cases}
d = 0 & \text{if } a \ne 2 \quad (\text{with } b,c \text{ arbitrary}), \\
b = c = d = 0 & \text{if } a = 2.
\end{cases}
\]

In short:  
- $e$ must be zero (no coupling from row 3 to column 1),  
- and the lower $2\times2$ block $\begin{pmatrix}2 & d\\0 & 2\end{pmatrix}$ must be diagonalizable, which forces $d = 0$,  
- unless $a = 2$, in which case the whole matrix must be diagonal (so $b = c = d = 0$).
\end{correctionbox}


% ----------- EXERCISE 9 -------------
\section{}
Given 
\[
A = \begin{pmatrix}
1 & 1 & 0 \\
-1 & 0 & 0 \\
2 & 0 & -1
\end{pmatrix}
\]
Use Cayley-Hamilton theorem to show \( A \) is invertible and compute its inverse.

\begin{correctionbox}
\textbf{Step 1: Characteristic polynomial.}  
Compute $ P(\lambda) = \det(A - \lambda I) $:
Expand along the third column (only one non-zero entry):
\[
P(\lambda) = (-1-\lambda) \cdot \det\begin{pmatrix} 1-\lambda & 1 \\ -1 & -\lambda \end{pmatrix}
= (-1-\lambda)\big[(1-\lambda)(-\lambda) - (-1)(1)\big].
\]
So the characteristic polynomial is  
\[
P(\lambda) = -\lambda^3 + 0\lambda^2 + 0\lambda + 1 = -\lambda^3 + 1,
\]
or equivalently (multiplying by $-1$, which doesn't affect Cayley–Hamilton):
\[
\chi_A(\lambda) = \lambda^3 - 1.
\]

\textbf{Step 2: Apply Cayley–Hamilton.}  
The theorem states $ \chi_A(A) = 0 $, so:
\[
A^3 - I = 0 \quad \Longrightarrow \quad A^3 = I.
\]

\textbf{Step 3: Invertibility and inverse.}  
From $ A^3 = I $, multiply both sides by $ A^{-1} $ (if it exists) to get $ A^2 = A^{-1} $. But more directly:  
\[
A \cdot A^2 = A^3 = I,
\]
so $ A $ is invertible and 
\[
A^{-1} = A^2.
\]

\textbf{Step 4: Compute $ A^2 $.}
\[
A^2 = A \cdot A = 
\begin{pmatrix}
1 & 1 & 0 \\
-1 & 0 & 0 \\
2 & 0 & -1
\end{pmatrix}
\begin{pmatrix}
1 & 1 & 0 \\
-1 & 0 & 0 \\
2 & 0 & -1
\end{pmatrix}
=
\begin{pmatrix}
0 & 1 & 0 \\
-1 & -1 & 0 \\
0 & 2 & 1
\end{pmatrix}.
\]

\textbf{Answer:} $ A $ is invertible and 
\[
A^{-1} = A^2 = \begin{pmatrix}
0 & 1 & 0 \\
-1 & -1 & 0 \\
0 & 2 & 1
\end{pmatrix}.
\]
\end{correctionbox}

% ----------- EXERCISE 10 -------------
\section{}
Let \( E \) be a vector space over \( K \) of dimension \( n \), and \( u \) an endomorphism of \( E \).\\
Show that for all \( p \in \mathbb{N} \), \( u^p \in Vect(I_d, u, \ldots, u^{m-1}) \).\\
Application: Let \( A \) be the matrix
\[
\begin{pmatrix}
2 & 0 & 4 \\
3 & -4 & 12 \\
1 & -2 & 5
\end{pmatrix}
\]
For \( p \in \mathbb{N} \), write \( A^p \) as a linear combination of \( I, A, A^2 \).

\begin{correctionbox}
Let $ E $ be a vector space over $ K $ with $ \dim E = n $, and $ u \in \mathrm{End}(E) $.  
Since $ \dim \mathrm{End}(E) = n^2 $, the $ n^2 + 1 $ endomorphisms $ I, u, u^2, \dots, u^{n^2} $ are linearly dependent. Hence there exists a non-zero polynomial $ P $ such that $ P(u) = 0 $. Let $ m = \deg(\mu_u) \leq n $ be the degree of the minimal polynomial of $ u $. Then
\[
u^m \in \mathrm{Vect}(I, u, \dots, u^{m-1}),
\]
and by induction, $ u^p \in \mathrm{Vect}(I, u, \dots, u^{m-1}) $ for all $ p \in \mathbb{N} $.

\medskip

\noindent\textbf{Application.} Let 
\[
A = \begin{pmatrix}
2 & 0 & 4 \\
3 & -4 & 12 \\
1 & -2 & 5
\end{pmatrix}.
\]

Compute the characteristic polynomial:
\[
\chi_A(\lambda) = \det(A - \lambda I) = -\lambda(\lambda - 1)(\lambda - 2).
\]

By the Cayley–Hamilton theorem, $ \chi_A(A) = 0 $, so
\[
A^3 - 3A^2 + 2A = 0 \quad \Longrightarrow \quad A^3 = 3A^2 - 2A.
\]

This recurrence allows us to reduce any higher power of $ A $ to a linear combination of $ I $, $ A $, and $ A^2 $. In fact, since the relation contains no $ I $ term, for $ p \geq 1 $ the powers lie in $ \mathrm{Vect}(A, A^2) $, but we include $ I $ for $ p = 0 $.

Thus, for every $ p \in \mathbb{N} $, there exist scalars $ \alpha_p, \beta_p, \gamma_p \in K $ such that
\[
A^p = \alpha_p I + \beta_p A + \gamma_p A^2,
\]
with the recurrence (for $ p \geq 3 $):
\[
A^p = 3A^{p-1} - 2A^{p-2}.
\]

\end{correctionbox}

\end{document}